\section{Large private convex quadratic blocks}\label{sec:recourse}

This section studies the recourse model of \cref{def:rec-model} with the
private-degree-two hierarchy of \cref{def:rec-hierarchy}. Each bag carries,
besides its shared coordinates, a private vector that may be much larger
than the bag, enters the objective as a convex quadratic, and is
constrained by linear inequalities. The hierarchy smooths only the shared
coordinates; private moments are kept at degree two at every order.

We prove four results. If the private polytopes do not depend on the
shared coordinates, the gap is $O(r^{-2})$ (\cref{thm:fixed-recourse}). The
redundant private quadratic bounds (R3) and private convexity cannot be
dropped from this statement (\cref{prop:private-bounds-needed,prop:convexity-needed}).
If the polytopes move affinely with the shared coordinates, have a fixed
constraint matrix, and are nonempty at every shared point, the gap is
$O(r^{-1})$ (\cref{thm:affine-recourse}), and this order is attained by a
three-variable example (\cref{thm:affine-sharp}). Finally, for every
feasible model the moment value of the hierarchy equals its certificate
value, and certificates exist at every level strictly below it
(\cref{prop:rec-dual}).

The mechanism is the one of \cref{sec:kernels}, applied to the shared
coordinates only. A common kernel turns each local functional into a
nonnegative density on its shared bag box, together with a positive
semidefinite matrix of conditional private moments. The densities have
exactly matching separator marginals and glue to a global shared law. The
private decision of each bag is then a measurable function of its own
shared coordinates: the conditional private mean in the fixed case, and a
Euclidean projection of that mean in the affine case. No private moment is
ever compared across bags, so no matrix-valued gluing is needed.

\subsection{Constants and kernel facts}\label{sec:rec-constants}

Throughout \cref{sec:recourse,sec:regularity}, $H_b(x)$ is the augmented
coefficient matrix of \eqref{eq:rec-objective}, a symmetric matrix of order
$p_b+1$ with polynomial entries, and
\begin{equation}\label{eq:rec-H-expansion}
  H_b(x)=\sum_{\alpha}H_{b\alpha}T_\alpha(x),\qquad
  H_{b\alpha}\in\R^{(p_b+1)\times(p_b+1)}\ \text{symmetric},
\end{equation}
is its tensor Chebyshev expansion in the shared variables of $B_b$. For a
matrix $M$ write $\norm{M}_{1}=\sum_{i,j}\abs{M_{ij}}$ (both off-diagonal
entries are counted). Put
\begin{equation}\label{eq:rec-budget}
  d=\max\{\abs{\alpha}:H_{b\alpha}\ne0\text{ for some }b\},\qquad
  \Abud_b=\sum_\alpha\norm{H_{b\alpha}}_{1}\sum_{i\in B_b}\alpha_i^2,\qquad
  \Abud=\sum_b\Abud_b,
\end{equation}
with $d=0$ for the zero objective. In \cref{sec:recourse,sec:regularity}
the symbol $\Abud$ always denotes this matrix budget. In terms of the
objective, $\Abud_b$ adds the weighted absolute Chebyshev coefficients of
$c_b$, of each entry of $q_b$ once, and of each ordered entry of $Q_b$.
When $p_b=0$ it equals $\Abud(f_b)$ of \eqref{eq:set-norms}. Coefficients
that do not depend on $x$ contribute nothing.

We use the kernel of \cref{lem:jackson}. For an integer $m\ge2$ let
$b_j=(m-\abs{j})_+$ for $j\in\Z$, $a_k=\sum_{j\in\Z}b_jb_{j-k}$, and
$\kmult_k=a_k/a_0$ for $0\le k\le2m-2$, $\kmult_k=0$ for $k>2m-2$. Then
\begin{equation}\label{eq:rec-circle-kernel}
  J_m(t)=\frac1{a_0}\Bigl\lvert\sum_{j=0}^{m-1}e^{ijt}\Bigr\rvert^4
  =1+2\sum_{k\ge1}\kmult_k\cos(kt)\ge0,\qquad
  \Kj(x,u)=1+2\sum_{k\ge1}\kmult_kT_k(x)T_k(u),
\end{equation}
and \cref{lem:jackson} gives $\Kj\ge0$ on $[-1,1]^2$,
$\int\Kj(x,u)\,d\arcs(u)=1$, $\int T_k(u)\Kj(x,u)\,d\arcs(u)=\kmult_kT_k(x)$
for all $k\ge0$, and $0\le1-\kmult_k\le3k^2/\Dm$, where $\Dm=2m^2+1$. For a
bag put
\[
  K_{m,b}(x,u)=\prod_{i\in B_b}\Kj(x_i,u_i),\qquad
  \Gamma_\alpha=\prod_{i\in B_b}\kmult_{\alpha_i},
\]
so that $\int T_\alpha(u)K_{m,b}(x,u)\,d\arcs^{B_b}(u)=\Gamma_\alpha
T_\alpha(x)$. Empty products equal one, and $\arcs^{\emptyset}$ is the unit
mass on the single point of $[-1,1]^{\emptyset}$; an empty bag therefore
causes no exception below. Define
\begin{equation}\label{eq:rec-Vm}
  V_m=\frac{3(4m-3)}{2m\Dm}=\frac6\Dm\cdot\frac{4m-3}{4m}<\frac6\Dm .
\end{equation}

\begin{lemma}[Kernel facts]\label{lem:rec-kernel}
Let $m\ge2$.
\begin{enumerate}[label=(\alph*)]
\item $a_0=m\Dm/3$, $1-\kmult_1=3/\Dm$, and
  $1-\kmult_2=(4m-3)/a_0$.
\item For every $x\in[-1,1]$,
  \begin{equation}\label{eq:rec-displacement-identity}
    \int\Kj(x,u)\,(x-u)\,d\arcs(u)=(1-\kmult_1)\,x,\qquad
    \int\Kj(x,u)\,(x-u)^2\,d\arcs(u)=V_m+\frac{3-2m}{a_0}\,x^2\le V_m .
  \end{equation}
\item For a bag $B_b$ and a fixed $u\in[-1,1]^{B_b}$ there are polynomials
  $p_{I,l}$ in $x_{B_b}$, finitely many for each $I\subseteq B_b$, with
  \begin{equation}\label{eq:rec-kernel-cert}
    K_{m,b}(x,u)=\sum_{I\subseteq B_b}w_I(x)\sum_lp_{I,l}(x)^2,\qquad
    \deg p_{I,l}\le v_b(m-1)-\abs{I}.
  \end{equation}
\end{enumerate}
\end{lemma}

\begin{proof}
(a) We have $a_0=\sum_{\abs{j}<m}(m-\abs{j})^2=m^2+2\sum_{i=1}^{m-1}i^2
=(2m^3+m)/3=m\Dm/3$. Since $\sum_jb_j^2=\sum_jb_{j-k}^2=a_0$,
\[
  a_0-a_k=\tfrac12\sum_{j\in\Z}(b_j-b_{j-k})^2 .
\]
For $k=1$ there are exactly $2m$ nonzero differences, each of absolute value
one, so $a_0-a_1=m$ and $1-\kmult_1=m/a_0=3/\Dm$. For $k=2$ the nonzero
differences $b_j-b_{j-2}$ are $1$ at $j=1-m$, $2$ for $2-m\le j\le0$, $-2$
for $2\le j\le m$, and $-1$ at $j=m+1$. Their squares sum to
$2+8(m-1)$, so $a_0-a_2=4m-3$.

(b) Integrate against $\Kj(x,u)\,d\arcs(u)$, using $u=T_1(u)$,
$u^2=\tfrac12(1+T_2(u))$ and the eigenvalue relation:
$\int\Kj(x,u)u\,d\arcs(u)=\kmult_1x$ and
$\int\Kj(x,u)u^2\,d\arcs(u)=\tfrac12+\tfrac12\kmult_2(2x^2-1)$. Hence the
first integral is $(1-\kmult_1)x$ and the second is
\[
  x^2-2\kmult_1x^2+\tfrac12+\tfrac12\kmult_2(2x^2-1)
  =\tfrac12(1-\kmult_2)+\bigl(2(1-\kmult_1)-(1-\kmult_2)\bigr)x^2 .
\]
By (a), $\tfrac12(1-\kmult_2)=(4m-3)/(2a_0)=V_m$ and
$2(1-\kmult_1)-(1-\kmult_2)=(2m-4m+3)/a_0=(3-2m)/a_0<0$.

(c) For fixed $u_i$, the polynomial $x_i\mapsto\Kj(x_i,u_i)$ is nonnegative
on $[-1,1]$ and has degree at most $2m-2$. By \cref{lem:interval-sos} it
equals $\sigma_{0,i}+(1-x_i^2)\sigma_{1,i}$ with $\sigma_{0,i}\in\sos_{2m-2}$
and $\sigma_{1,i}\in\sos_{2m-4}$. Multiply these representations over
$i\in B_b$ and expand. The term indexed by $I\subseteq B_b$ is $w_I$ times
a product of sums of squares, which is a sum of squares of polynomials of
degree at most $(v_b-\abs{I})(m-1)+\abs{I}(m-2)=v_b(m-1)-\abs{I}$.
\end{proof}

The representation in (c) need not depend continuously on $u$; it is used
at each fixed $u$ separately.

\begin{remark}[Private and shared dimensions]\label{rem:rec-dimension}
The two kinds of variables enter the hierarchy differently. The order $r$
needed for a given accuracy is set by the kernel, which acts only on shared
coordinates, through $m=\lfloor r/w\rfloor+1$ or $m=\lceil r/w\rceil$ below.
For fixed $r$, by \eqref{eq:rec-size}, the shared bag size $v_b$ appears
in the exponent of the block sizes, about $r^{v_b}$, while the private
dimension enters only through the factor $p_b+1$ of the largest block, the
number $J_b+p_b$ of scalar localizers, and at most $\binom{p_b+2}{2}$ private
monomials. Accuracy constants are a different matter: $\Abud$, and in the
affine case the Hoffman and gradient constants of \cref{thm:affine-recourse},
can grow with $p_b$ and with the data. These statements concern the sizes
of finite semidefinite programs. We make no claim about the time needed to
solve them, about their numerical conditioning, or about a practical
advantage over other formulations; see also \cref{rem:grid-baseline}.
\end{remark}

\subsection{Conditional moment matrices}\label{sec:rec-conditional}

Fix a feasible family $(L_b)$ of order $r$ for \cref{def:rec-hierarchy}.
Restricted to polynomials in $x_{B_b}$ alone, (R1) with $s=0$ shows that
$L_b$ is $\Pre_r(B_b)$-positive, so \cref{lem:cheb-moment,lem:moment-cs}
apply to it. We write $\zy=(z_0,z_1,\dots,z_{p_b})^\top$ with $z_0=1$ and
$z_k=y_{b,k}$.

\begin{lemma}[Mixed moments]\label{lem:mixed-moment}
For every multi-index $\alpha$ with $\abs{\alpha}\le r$ and all
$0\le i,j\le p_b$,
\[
  \abs{L_b(T_\alpha z_iz_j)}\le1 .
\]
\end{lemma}

\begin{proof}
By \cref{lem:cheb-moment}, $0\le L_b(T_\alpha^2)\le1$. For $k\ge1$, (R1)
gives $L_b(T_\alpha^2y_{b,k}^2)\ge0$, and (R3) with $I=\emptyset$ and
$s=T_\alpha$ gives $L_b(T_\alpha^2(1-y_{b,k}^2))\ge0$. Hence
$0\le L_b(T_\alpha^2z_i^2)\le1$ for every $i$, and in particular
$0\le L_b(z_j^2)\le1$. Apply \cref{lem:moment-cs} with $G=1$ to the space of
polynomials $s_0(x)+s(x)^\top y_b$ with $\deg s_k\le r$, on which (R1) with
$I=\emptyset$ is a positive semidefinite form; both $T_\alpha z_i$ and $z_j$
belong to it. Thus $L_b(T_\alpha z_iz_j)^2\le L_b(T_\alpha^2z_i^2)L_b(z_j^2)
\le1$.
\end{proof}

The private quadratic bounds (R3) are used here: affine localizers alone
do not bound second private moments. With no shared variables, the
functional $L(1)=1$, $L(y)=0$, $L(y^2)=R$ on $\Rle{y}{2}$ satisfies
(R1)--(R2) for $P=[-1,1]$ and every $R\ge0$.

For a kernel degree $m\ge2$ and a bag $b$, and for every output point
$u\in[-1,1]^{B_b}$, define the conditional moments
\begin{equation}\label{eq:rec-conditional}
  h_b(u)=L_b\bigl(K_{m,b}(\cdot,u)\bigr),\quad
  \psi_b(u)=L_b\bigl(K_{m,b}(\cdot,u)\,y_b\bigr),\quad
  Y_b(u)=L_b\bigl(K_{m,b}(\cdot,u)\,y_by_b^\top\bigr),
\end{equation}
where $L_b$ acts on $(x,y_b)$ entrywise and $u$ is a parameter, and set
\[
  M_b(u)=L_b\bigl(K_{m,b}(\cdot,u)\,\zy\zy^\top\bigr)
  =\begin{pmatrix}h_b(u)&\psi_b(u)^\top\\ \psi_b(u)&Y_b(u)\end{pmatrix}.
\]
These are defined when $2v_b(m-1)\le2r$, and they are polynomials in $u$.
When $2v_b(m-1)\le2r-2$ we also use
\[
  U_b(u)=L_b\bigl(K_{m,b}(\cdot,u)\,x_{B_b}\bigr)\in\R^{v_b},\qquad
  \widehat M_b(u)=L_b\bigl(K_{m,b}(\cdot,u)\,\zeta\zeta^\top\bigr),\quad
  \zeta=(1,x_{B_b}^\top,y_b^\top)^\top .
\]

\begin{lemma}[Conditional matrices]\label{lem:conditional-matrix}
Let $m\ge2$ and $\tau\in\{0,1\}$ satisfy $v_b(m-1)\le r-\tau$. Then for every
$u\in[-1,1]^{B_b}$:
\begin{enumerate}[label=(\alph*)]
\item $M_b(u)\succeq0$ and $0\le Y_b(u)_{kk}\le h_b(u)$ for every $k$. If
  $h_b(u)=0$ then $M_b(u)=0$.
\item For every fiber row $j$ with $\epsilon_{bj}\le\tau$, that is, every
  row if $\tau=1$ and every row with $E_b$-row zero if $\tau=0$,
  \[
    a_{bj}h_b(u)+(E_bU_b(u))_j-(A_b\psi_b(u))_j\ge0,
  \]
  where the middle term is absent when $\tau=0$ (it is then multiplied by
  a zero row).
\item If $\tau=1$, then $\widehat M_b(u)\succeq0$ and
  $0\le L_b(K_{m,b}(\cdot,u)x_i^2)\le h_b(u)$ for $i\in B_b$. If moreover
  $h_b(u)=0$, then $\widehat M_b(u)=0$; in particular $U_b(u)=0$.
\item If $h_b(u)>0$, then $\bar y_b(u)=\psi_b(u)/h_b(u)\in[-1,1]^{p_b}$, and
  $A_b\bar y_b(u)\le a_b$ holds in every row with zero $E_b$-row. If
  $\tau=1$, then also $\bar x_b(u)=U_b(u)/h_b(u)\in[-1,1]^{B_b}$ and
  \begin{equation}\label{eq:rec-source-feasible}
    \bar y_b(u)\in P_b\bigl(\bar x_b(u)\bigr).
  \end{equation}
\end{enumerate}
\end{lemma}

\begin{proof}
Write $K=K_{m,b}(\cdot,u)$ and use \eqref{eq:rec-kernel-cert}; each square
polynomial $p_{I,l}$ has degree at most $r-\tau-\abs{I}$.

(a) For a constant vector $c\in\R^{p_b+1}$,
$c^\top M_b(u)c=\sum_{I,l}L_b\bigl(w_I(p_{I,l}\,c^\top\zy)^2\bigr)$. Each
$p_{I,l}\,c^\top\zy$ has the form $s_0+s^\top y_b$ with all $s_k$ of degree
at most $r-\abs{I}$, so (R1) gives $c^\top M_b(u)c\ge0$. Similarly
$h_b(u)-Y_b(u)_{kk}=\sum_{I,l}L_b(w_Ip_{I,l}^2(1-y_{b,k}^2))\ge0$ by (R3),
and $Y_b(u)_{kk}\ge0$ because $M_b(u)\succeq0$. If $h_b(u)=0$, every
diagonal entry of $M_b(u)$ vanishes, and so does every entry, because
$M_{ij}^2\le M_{ii}M_{jj}$ for a positive semidefinite matrix.

(b) The left side equals $L_b(K\,\ell_{bj})=\sum_{I,l}L_b(w_Ip_{I,l}^2
\ell_{bj})$. Since $\deg p_{I,l}\le r-\tau-\abs{I}\le r-\abs{I}-\epsilon_{bj}$,
each term is nonnegative by (R2).

(c) For $c=(c_0,c_x,c_y)$, $c^\top\widehat M_b(u)c=\sum_{I,l}L_b\bigl(w_I
(p_{I,l}(c_0+c_x^\top x)+p_{I,l}c_y^\top y_b)^2\bigr)$. The coefficient
polynomials have degree at most $r-1-\abs{I}+1=r-\abs{I}$, so (R1) applies.
Next, $h_b(u)-L_b(Kx_i^2)=\sum_{I,l}L_b(w_I(1-x_i^2)p_{I,l}^2)$. If
$i\notin I$, the term is $w_{I\cup\{i\}}p_{I,l}^2$ with
$\deg p_{I,l}\le r-\abs{I\cup\{i\}}$. If $i\in I$, it is
$w_{I\setminus\{i\}}\bigl((1-x_i^2)p_{I,l}\bigr)^2$, whose square polynomial
has degree at most $r-1-\abs{I}+2=r-\abs{I\setminus\{i\}}$. Both are
allowed in (R1), so $L_b(Kx_i^2)\le h_b(u)$. Nonnegativity of the diagonal
of $\widehat M_b(u)$ and the minor argument of (a) finish the proof, the
private diagonal being bounded by (a).

(d) By the $2\times2$ minors of $M_b(u)$ and (a),
$\psi_b(u)_k^2\le h_b(u)Y_b(u)_{kk}\le h_b(u)^2$, so $\abs{\bar y_{b,k}}\le1$.
The same argument with (c) gives $\abs{\bar x_{b,i}}\le1$ when $\tau=1$.
Dividing the inequalities of (b) by $h_b(u)>0$ gives the remaining
claims; when $\tau=1$ every row of $A_b\bar y_b\le a_b+E_b\bar x_b$ holds,
and together with $\bar y_b\in[-1,1]^{p_b}$ this is
\eqref{eq:rec-source-feasible}.
\end{proof}

Part (d) guarantees feasibility at the conditional mean $\bar x_b(u)$ of the
\emph{input} shared variables. In the affine case the private mean need not
be feasible at the output $u$. Bounding the cost of repairing it there gives
the general $O(r^{-1})$ term in \cref{thm:affine-recourse}.

\begin{lemma}[Matrix cost]\label{lem:rec-matrix-cost}
Let $m\ge2$, $v_b(m-1)\le r$, and $d\le r$. Then
\[
  \int\ip{H_b(u)}{M_b(u)}\,d\arcs^{B_b}(u)\le L_b(f_b)+\frac{3\Abud_b}{\Dm},
\]
where $\ip{H}{M}=\sum_{i,j}H_{ij}M_{ij}$.
\end{lemma}

\begin{proof}
By \eqref{eq:rec-H-expansion}, $\ip{H_b(u)}{M_b(u)}=\sum_\alpha\sum_{i,j}
(H_{b\alpha})_{ij}T_\alpha(u)M_b(u)_{ij}$. Expand $K_{m,b}(x,u)$ as a finite
sum of products of polynomials in $x$ and in $u$. By linearity of $L_b$,
\[
  \int T_\alpha(u)M_b(u)_{ij}\,d\arcs^{B_b}(u)
  =L_b\Bigl(z_iz_j\int T_\alpha(u)K_{m,b}(x,u)\,d\arcs^{B_b}(u)\Bigr)
  =\Gamma_\alpha L_b(T_\alpha z_iz_j).
\]
Also $L_b(f_b)=\sum_\alpha\sum_{i,j}(H_{b\alpha})_{ij}L_b(T_\alpha z_iz_j)$. Every
$\alpha$ with $H_{b\alpha}\ne0$ has $\abs{\alpha}\le d\le r$, so
\cref{lem:mixed-moment} bounds the difference by
$\sum_\alpha\norm{H_{b\alpha}}_1(1-\Gamma_\alpha)$. Since all $\kmult_k$ lie in
$[0,1]$, induction on the number of factors gives
$0\le1-\Gamma_\alpha\le\sum_i(1-\kmult_{\alpha_i})\le3\sum_i\alpha_i^2/\Dm$.
Multipliers with $\alpha_i>2m-2$ are zero; this case is included.
\end{proof}

\begin{lemma}[Consistent shared laws]\label{lem:rec-glue}
Let $m\ge2$ with $v_b(m-1)\le r$ for every $b$. Then each $h_b$ is
nonnegative on $[-1,1]^{B_b}$, $\int h_b\,d\arcs^{B_b}=1$, and for every edge
$e=bc$ the laws $h_b\,d\arcs^{B_b}$ and $h_c\,d\arcs^{B_c}$ have the same
marginal on $[-1,1]^{S_e}$. Let $\phi_b:[-1,1]^{B_b}\to\R^{p_b}$ be Borel maps
with $\phi_b(u)\in P_b(u)$ for every $u$. Then there is a probability law
$\nu$ on the feasible set of the recourse model with
\[
  \int f\,d\nu=\sum_b\int f_b\bigl(u,\phi_b(u)\bigr)h_b(u)\,d\arcs^{B_b}(u),
\]
and consequently $\fstar$ is at most the right side.
\end{lemma}

\begin{proof}
Nonnegativity is $M_b(u)_{00}\ge0$ in \cref{lem:conditional-matrix}(a).
Integrating $\Kj(x_i,u_i)$ in $u_i$ gives one, so
$\int h_b\,d\arcs^{B_b}=L_b(1)=1$, and the marginal of $h_b\,d\arcs^{B_b}$ on
$S_e$ has density $L_b\bigl(\prod_{i\in S_e}\Kj(x_i,u_i)\bigr)$ with respect to
$\arcs^{S_e}$. The argument of $L_b$ is a polynomial in $x_{S_e}$ of degree at
most $2v_b(m-1)\le2r$, so (R4) makes this density the same for $b$ and $c$.
By \cref{lem:gluing} there is a law $\nu^{\mathrm{sh}}$ on $[-1,1]^n$ with
marginals $h_b\,d\arcs^{B_b}$. Let $\nu$ be its image under
$x\mapsto(x,(\phi_b(x_{B_b}))_b)$. It is concentrated on the feasible set,
and its objective is the displayed sum because each $f_b$ depends only on
$(x_{B_b},y_b)$. The objective is at least $\fstar$ at every feasible
point.
\end{proof}

\subsection{Finiteness, attainment, and certificates}\label{sec:rec-dual}

For the box hierarchies, product arcsine moments are strictly feasible and
\cref{lem:slater} gives zero duality gap and dual attainment. That argument
does not transfer to the recourse hierarchy: a fiber may be
lower-dimensional, for instance when an equality is written as two
inequalities, and then no feasible family has positive definite localizing
matrices. We use instead an order-unit argument, which needs no strict
feasibility and gives slightly less: certificates at every level below the
moment value, without a claim at the moment value itself.

Let $V_{b,r}=\Rle{x_{B_b}}{2r}\otimes\Rle{y_b}{2}$, let $C_{b,r}\subseteq
V_{b,r}$ be the convex cone of finite sums of the polynomials tested in
(R1)--(R3), with the degree allowances of \cref{def:rec-hierarchy}, and let
\[
  W_r=\sum_bV_{b,r},\qquad C_r=\sum_bC_{b,r}\subseteq W_r
\]
be the sums inside the global polynomial ring in $(x,y_1,\dots,y_t)$. The
\emph{certificate value} of the private-degree-two hierarchy is
\begin{equation}\label{eq:rec-certificate}
  \lambda^{\mathrm{rec}}_r=\sup\bigl\{\lambda\in\R:\ f-\lambda\in C_r\bigr\},
  \qquad f=\sum_bf_b ,
\end{equation}
that is, the supremum of $\lambda$ for which $f-\lambda=\sum_b\sigma_b$ as
polynomials, with $\sigma_b\in C_{b,r}$.

\begin{lemma}[Decompositions of zero with private variables]
\label{lem:rec-zero}
If $p_b\in V_{b,r}$ and $\sum_bp_b=0$, then there are polynomials
$\phi_e\in\Rle{x_{S_e}}{2r}$ such that each $p_b$ is the signed sum of the
$\phi_e$ over the edges incident to $b$, as in \cref{lem:zero-decomposition}.
In particular, every feasible family $(L_b)$ defines a linear functional
$F$ on $W_r$ by $F(\sum_bp_b)=\sum_bL_b(p_b)$.
\end{lemma}

\begin{proof}
A monomial containing a variable of $y_b$ occurs only in $p_b$, because
private vectors of different bags are distinct variables. Its coefficient
in $\sum_bp_b=0$ is therefore its coefficient in $p_b$, which must vanish.
Hence each $p_b$ is a polynomial in $x_{B_b}$ of degree at most $2r$, and
\cref{lem:zero-decomposition} applies. If $\sum_bp_b=\sum_bp'_b$, the
differences $p_b-p'_b$ are signed sums of separator polynomials of degree at
most $2r$, on which adjacent functionals agree by (R4); so $F$ is well
defined.
\end{proof}

\begin{lemma}[Order unit]\label{lem:rec-order-unit}
For every bag $b$, every multi-index $\alpha$ with $\abs{\alpha}\le2r$, and
all $0\le i,j\le p_b$, both $1+x^\alpha z_iz_j$ and $1-x^\alpha z_iz_j$
belong to $C_{b,r}$.
\end{lemma}

\begin{proof}
Write $\alpha=\beta+\beta'$ with $\abs\beta,\abs{\beta'}\le r$, and put
$g=x^\beta z_i$ and $g'=x^{\beta'}z_j$. Order the coordinates of $B_b$. For
$\abs{\beta}\le r$, telescoping and $1-x_h^{2c}=(1-x_h^2)\sum_{l=0}^{c-1}
x_h^{2l}$ give
\begin{equation}\label{eq:rec-telescope-monomial}
  1-x^{2\beta}=\sum_{h\in B_b}\sum_{l=0}^{\beta_h-1}(1-x_h^2)
  \Bigl(x_h^l\prod_{h'<h}x_{h'}^{\beta_{h'}}\Bigr)^2 .
\end{equation}
Each square polynomial has degree at most $\abs\beta-1\le r-1$, so every
term is allowed in (R1) with $I=\{h\}$. If $i=0$, this shows $1-g^2\in
C_{b,r}$. If $i\ge1$, write $1-g^2=(1-x^{2\beta})+x^{2\beta}(1-y_{b,i}^2)$;
the last term is allowed in (R3) with $I=\emptyset$ and $s=x^\beta$. The
same holds for $g'$. Finally
\[
  1\pm gg'=\tfrac12\bigl[(1-g^2)+(1-g'^2)+(g\pm g')^2\bigr],
\]
and $(g\pm g')^2$ is allowed in (R1) with $I=\emptyset$, because $g\pm g'$
is affine in $y_b$ with shared coefficients of degree at most $r$.
\end{proof}

\begin{proposition}[Moment and certificate values]\label{prop:rec-dual}
Let $r\ge1$ with $f_b\in V_{b,r}$ for every bag, and suppose the recourse
model has at least one feasible point
(this holds under nonempty fixed polytopes or complete recourse). Then:
\begin{enumerate}[label=(\roman*)]
\item the set of feasible families of order $r$ is compact, every local
  moment of a feasible family has absolute value at most one, and
  $\rhorec_r$ is finite and attained;
\item $\lambda^{\mathrm{rec}}_r=\rhorec_r$;
\item $f-\lambda\in C_r$ for every $\lambda<\rhorec_r$;
\item if $\fstar-\rhorec_r\le\mathcal E$, then
  $f-(\fstar-\mathcal E-\varepsilon)\in C_r$ for every $\varepsilon>0$.
\end{enumerate}
We make no claim that $f-\rhorec_r\in C_r$.
\end{proposition}

\begin{proof}
\emph{Global functionals.} By \cref{lem:rec-zero}, a feasible family
defines a linear functional $F$ on $W_r$ with $F(1)=1$, $F\ge0$ on $C_r$, and
$F(f)=\sum_bL_b(f_b)$. Conversely, if $F$ is a linear functional on $W_r$
with $F(1)=1$ and $F\ge0$ on $C_r$, its restrictions $L_b=F|_{V_{b,r}}$
satisfy (R0)--(R3), and they satisfy (R4) because a separator polynomial lies
in both $V_{b,r}$ and $V_{c,r}$. Hence
\[
  \rhorec_r=\inf\bigl\{F(f):\ F\in C_r^*,\ F(1)=1\bigr\},
\]
where $C_r^*$ is the dual cone in the dual space of $W_r$.

\emph{Interior point.} Choose a basis of $W_r$ consisting of monomials. Each
basis monomial lies in some $V_{b,r}$ and has the form $x^\alpha z_iz_j$ of
\cref{lem:rec-order-unit}, so $1\pm e\in C_r$ for every basis element $e$.
If $g=\sum_ec_ee$ with $\sum_e\abs{c_e}\le\varepsilon$, then
\[
  \varepsilon+g=\Bigl(\varepsilon-\sum_e\abs{c_e}\Bigr)\cdot1
  +\sum_e\abs{c_e}\bigl(1+\operatorname{sign}(c_e)\,e\bigr)\in C_r,
\]
using $1\in C_r$. Hence $C_r$ contains a coefficient ball around each
$\varepsilon\cdot1$, $\varepsilon>0$; in particular $1$ is an interior point of
$C_r$ in $W_r$. We do not use, and do not claim, that $C_r$ is closed.

(i) If $F\in C_r^*$, then $\abs{F(e)}\le F(1)$ for every basis monomial $e$.
Thus every normalized $F\in C_r^*$ has all coordinates in $[-1,1]$, and the
same bound holds for the local moments of every feasible family. The set of
such $F$ is closed, being cut out by the closed half-spaces $F(c)\ge0$ and
the hyperplane $F(1)=1$. It is nonempty, because evaluation at a feasible
point $(x,y)$ is nonnegative on every generator: $w_I\ge0$ on the box, and
$\ell_{bj}\ge0$ and $1-y_{b,k}^2\ge0$ at feasible points. A linear function
attains its minimum on a nonempty compact set.

(iii) Let $\lambda<\rhorec_r$ and suppose $p=f-\lambda\notin C_r$. Since $C_r$
is convex with nonempty interior, the separation theorem for convex sets in
finite dimension~\cite[Theorem~2.4.2]{bental2001-lectures-modern-convex} gives a nonzero
linear functional $F$ on $W_r$ with $F(c)\ge F(p)$ for all $c\in C_r$.
Because $C_r$ is a cone containing $0$, this gives $F\in C_r^*$ and
$F(p)\le0$. By (i), $F\ne0$ forces $F(1)>0$. Then $G=F/F(1)$ is normalized
and $G(f)\le\lambda<\rhorec_r$, contradicting the formula for $\rhorec_r$.

(ii) If $f-\lambda\in C_r$, then every normalized $F\in C_r^*$ has
$F(f)\ge\lambda$, so $\lambda^{\mathrm{rec}}_r\le\rhorec_r$; with (iii),
equality holds.

(iv) Apply (iii) with $\lambda=\fstar-\mathcal E-\varepsilon<\rhorec_r$.
\end{proof}

Each element of $C_r$ is a finite sum of the tested polynomials, so (iii)
asserts a finite certificate with real coefficients: positive semidefinite
Gram matrices for (R1) and sum-of-squares multipliers for the scalar localizers.
The proposition gives no bound on the size of these coefficients, no
rational certificate, and no method to compute one. The order-unit
argument uses only (R1) and (R3). Without (R3) it fails, and
\cref{prop:private-bounds-needed} shows that $\rhorec_r$ can then be
$-\infty$. Every gap bound proved below for $\rhorec_r$ is also a bound on
$\fstar-\lambda^{\mathrm{rec}}_r$, by (ii).

\subsection{Fixed private polytopes}\label{sec:rec-fixed}

\begin{theorem}[Fixed private domain]\label{thm:fixed-recourse}
Consider the recourse model with a fixed private domain ($E_b=0$ for all
$b$), nonempty polytopes $P_b$, and private convexity. Let $r\ge\max(w,d)$
and $m=\lfloor r/w\rfloor+1$. Every feasible family $(L_b)$ of order $r$
admits a probability law $\nu$ on the feasible set with
\begin{equation}\label{eq:rec-fixed-rounding}
  \int f\,d\nu\le\sum_bL_b(f_b)+\frac{3\Abud}{\Dm}.
\end{equation}
Consequently
\[
  0\le\fstar-\rhorec_r=\fstar-\lambda^{\mathrm{rec}}_r
  \le\frac{3\Abud}{2m^2+1}\le\frac{3w^2\Abud}{2r^2}.
\]
In particular, if no entry of any $H_b$ depends on $x$, the hierarchy is
exact at every order $r\ge w$.
\end{theorem}

\begin{proof}
Here $v_b(m-1)\le w\lfloor r/w\rfloor\le r$, so \cref{lem:conditional-matrix}
applies with $\tau=0$; since $E_b=0$, every row has $\epsilon_{bj}=0$. Fix a
point $y_b^\circ\in P_b$ and define $\phi_b(u)=\bar y_b(u)$ if $h_b(u)>0$ and
$\phi_b(u)=y_b^\circ$ otherwise. The set $\{h_b>0\}$ is open and $\bar y_b$
is continuous there, so $\phi_b$ is Borel. By
\cref{lem:conditional-matrix}(d), $\phi_b(u)\in P_b$ for every $u$.

We claim that for every $u$,
\begin{equation}\label{eq:rec-jensen}
  h_b(u)\,f_b\bigl(u,\phi_b(u)\bigr)\le\ip{H_b(u)}{M_b(u)} .
\end{equation}
If $h_b(u)=0$, both sides vanish by \cref{lem:conditional-matrix}(a). If
$h=h_b(u)>0$, write $\psi=\psi_b(u)$, $Y=Y_b(u)$, and $\bar y=\psi/h$. Then
\[
  \ip{H_b(u)}{M_b(u)}-h\,f_b(u,\bar y)
  =\tr\bigl(Q_b(u)(Y-h^{-1}\psi\psi^\top)\bigr)\ge0,
\]
because $Y-h^{-1}\psi\psi^\top$ is the Schur complement of $h$ in
$M_b(u)\succeq0$ and $Q_b(u)\succeq0$.

Integrate \eqref{eq:rec-jensen}, apply \cref{lem:rec-matrix-cost} (its
hypotheses hold because $d\le r$), and sum over bags. \Cref{lem:rec-glue}
gives a law $\nu$ on the feasible set with
$\int f\,d\nu=\sum_b\int f_b(u,\phi_b(u))h_b(u)\,d\arcs^{B_b}\le
\sum_bL_b(f_b)+3\Abud/\Dm$. Taking the infimum over feasible families gives
the gap bound, and \cref{prop:rec-dual} gives the equality with
$\lambda^{\mathrm{rec}}_r$. Finally $m>r/w$ gives $2m^2+1>2r^2/w^2$.
\end{proof}

The comparison \eqref{eq:rec-fixed-rounding} is one-sided. The Jensen
step \eqref{eq:rec-jensen} can strictly decrease the private cost, so the
rounded objective can lie well below $\sum_bL_b(f_b)$; no bound on
$\abs{\int f\,d\nu-\sum_bL_b(f_b)}$ is claimed. The conditional-mean inequality
\eqref{eq:rec-jensen} is the quadratic case of the pseudomoment Jensen
inequality for SOS-convex polynomials~\cite{lasserre2009-convexity-sdp}.
Kahl and Henrion's partial lifting linearizes only a limited nonlinear
subset; it does not guarantee asymptotic convergence in
general~\cite{kahl2005-globally-optimal-estimates}.
Guo and Wang hold an SOS-convex decision-variable module at fixed degree
while increasing the uncertainty-variable order in a robust polynomial
matrix inequality hierarchy~\cite{wang2025-a-moment-sum-of-squares}.
Neither result supplies the sparse private-degree-two recourse rate proved here.
What the theorem adds is the quantitative rate for this sparse hierarchy,
with fixed private polytopes that need not be full-dimensional. Because the
private matrices are eliminated inside each bag and only scalar shared
densities are glued, the failure of generic sparse matrix hierarchies under
running intersection~\cite{miller2025-sparse-matrix} does not arise here.

Both hypotheses beyond the fixed domain are needed, in a strong sense.

\begin{proposition}[The private quadratic bounds are needed]
\label{prop:private-bounds-needed}
Remove (R3) from \cref{def:rec-hierarchy}. Take one bag with shared
coordinates $x=(x_1,x_2,x_3)$, one private variable $y$ with the two fiber
rows $1-y\ge0$ and $1+y\ge0$ (so $P=[-1,1]$ and $E=0$), and
\[
  f(x,y)=q(x)\,y^2,\qquad
  q(x)=x_1^4x_2^2+x_1^2x_2^4+x_3^6-3x_1^2x_2^2x_3^2 .
\]
This instance satisfies private convexity and has $\fstar=0$. For every
order $r\ge3$, the relaxation without (R3) has infimum $-\infty$.
\end{proposition}

The polynomial $q$ is the Motzkin polynomial. The proof, in
\cref{sec:app-rec-bounds}, shows that $q$ lies in no truncated preordering
$\Pre_r(\{1,2,3\})$, separates it from that closed cone by a normalized
functional $\mathcal L$, and uses the feasible functionals
$L_R(a(x)+b(x)y+c(x)y^2)=a(0)+R\,\mathcal L(c)$, whose objective $R\,\mathcal L(q)$
tends to $-\infty$.

\begin{proposition}[Private convexity is needed]\label{prop:convexity-needed}
Take one bag with one shared coordinate $x_1$, which does not appear in the
data, a private vector $y\in\R^2$ with fixed polytope
$P=\{y\in[-1,1]^2:\ y_1\ge0,\ y_2\ge0,\ y_1+y_2\le1\}$, and the objective
$f=-(y_1+y_2)^2$, which violates private convexity. Then $\fstar=-1$, while
$\rhorec_r\le-4$ for every $r\ge1$. Since $\Abud=0$ for this instance,
\cref{thm:fixed-recourse} fails without private convexity.
\end{proposition}

\begin{proof}
On $P$ we have $0\le y_1+y_2\le1$, so $\fstar=-1$. Let $\Lambda$ be the
linear functional on $\Rle{y}{2}$ with $\Lambda(1)=1$,
$\Lambda(y_1)=\Lambda(y_2)=\tfrac12$, and $\Lambda(y_1^2)=\Lambda(y_2^2)=
\Lambda(y_1y_2)=1$. Its moment matrix
\[
  \begin{pmatrix}1&\tfrac12&\tfrac12\\ \tfrac12&1&1\\ \tfrac12&1&1\end{pmatrix}
\]
is positive semidefinite: the Schur complement of the leading entry is
$\tfrac34\bigl(\begin{smallmatrix}1&1\\1&1\end{smallmatrix}\bigr)$. Define
$L(p)=\Lambda\bigl(p(0,\cdot)\bigr)$ on $\Rle{x_1}{2r}\otimes\Rle{y}{2}$,
evaluating the shared variable at zero. For a polynomial tested in (R1),
$p(0,y)=w_I(0)(s_0(0)+s(0)^\top y)^2$, and $\Lambda$ of a square of an affine
form is nonnegative. For (R2), $p(0,y)=w_I(0)s(0)^2\ell_j(y)$ with
$\Lambda(y_1)=\Lambda(y_2)=\tfrac12$ and $\Lambda(1-y_1-y_2)=0$. For (R3),
$\Lambda(1-y_k^2)=0$. There are no separator conditions. Hence $L$ is
feasible at every order, and
$L(f)=-\Lambda(y_1^2+2y_1y_2+y_2^2)=-4$.
\end{proof}

The shared coordinate in \cref{prop:convexity-needed} only places the
instance in the setting $w\ge1$ of \cref{def:junction-tree}; it plays no role.
The private moment matrix above has no representing measure on $P$, since
$\Lambda(y_1^2)=1>\Lambda(y_1)$ is impossible for a law on $[0,1]$.

\subsection{Affine recourse}\label{sec:rec-affine}

When the fibers move with $x$, the conditional private mean is guaranteed
to be feasible at the conditional input mean $\bar x_b(u)$ of
\cref{lem:conditional-matrix}(d), but need not be feasible at the output
$u$. We repair it by Euclidean projection onto the output fiber and bound
the cost using a Hoffman error bound. Write the fiber with its box rows as
\begin{equation}\label{eq:rec-stacked}
  P_b(x)=\bigl\{y:\ \widehat A_by\le\widehat a_b+\widehat E_bx\bigr\},\qquad
  \widehat A_b=\begin{pmatrix}A_b\\ I\\ -I\end{pmatrix},\quad
  \widehat a_b=\begin{pmatrix}a_b\\ \mathbf 1\\ \mathbf 1\end{pmatrix},\quad
  \widehat E_b=\begin{pmatrix}E_b\\ 0\\ 0\end{pmatrix}.
\end{equation}
By Hoffman's theorem~\cite{hoffman1952-approximate-solutions}, in the form
that is uniform over consistent right-hand
sides~\cite{pena2018-hoffman-constants}, there is a constant
$\kappa_b<\infty$, depending only on $\widehat A_b$, such that for every
vector $c$ for which $\{z:\widehat A_bz\le c\}$ is nonempty and every
$y\in\R^{p_b}$,
\begin{equation}\label{eq:rec-hoffman}
  \dist\bigl(y,\{z:\widehat A_bz\le c\}\bigr)\le
  \kappa_b\,\norm{(\widehat A_by-c)_+}_2 .
\end{equation}
Distances and norms of vectors are Euclidean, $\norm{E_b}_2$ is the
spectral norm, and $(\cdot)_+$ acts componentwise. Put
\begin{equation}\label{eq:rec-affine-constants}
  \Lambda_b=\max_{x\in[-1,1]^{B_b},\ y\in[-1,1]^{p_b}}
  \norm{q_b(x)+2Q_b(x)y}_2,\qquad
  \Gamma_b=\Lambda_b\kappa_b\norm{E_b}_2,\qquad
  \mathcal R=\sum_b\Gamma_b\sqrt{v_b}.
\end{equation}
Since $\nabla_yf_b=q_b+2Q_by$, $\Lambda_b$ is a Lipschitz constant of
$f_b(x,\cdot)$ on the private box, uniformly in $x$. A crude bound is
$\Lambda_b\le\sum_\alpha\norm{q_{b\alpha}}_2+2\sqrt{p_b}\sum_\alpha
\norm{Q_{b\alpha}}_2$, from the Chebyshev expansions of $q_b$ and $Q_b$.

\begin{lemma}[Continuous fibers and projections]\label{lem:rec-hoffman}
Assume complete recourse. For $x,x'\in[-1,1]^{B_b}$ the Hausdorff distance
satisfies $d_H(P_b(x),P_b(x'))\le\kappa_b\norm{E_b}_2\norm{x-x'}_2$, and the
map $(x,z)\mapsto\proj_{P_b(x)}(z)$ is continuous on
$[-1,1]^{B_b}\times\R^{p_b}$. Moreover, if $\bar x\in[-1,1]^{B_b}$,
$\bar y\in P_b(\bar x)$ and $u\in[-1,1]^{B_b}$, then
\begin{equation}\label{eq:rec-repair}
  \norm{\proj_{P_b(u)}(\bar y)-\bar y}_2\le\kappa_b\norm{E_b}_2
  \norm{\bar x-u}_2 .
\end{equation}
\end{lemma}

\begin{proof}
For \eqref{eq:rec-repair}: $P_b(u)$ is nonempty, compact, and convex, so the
projection exists and is unique. The box rows of $\widehat A_b\bar y-
\widehat a_b-\widehat E_bu$ are nonpositive, and in the rows of $A_b$,
$A_b\bar y-a_b-E_bu\le E_b(\bar x-u)$ because $A_b\bar y\le a_b+E_b\bar x$.
Hence $\norm{(\widehat A_b\bar y-\widehat a_b-\widehat E_bu)_+}_2\le
\norm{E_b(\bar x-u)}_2$, and \eqref{eq:rec-hoffman} gives
\eqref{eq:rec-repair}. Applying \eqref{eq:rec-repair} to every point of
$P_b(x')$ with $u=x$, and symmetrically, gives the Hausdorff bound.

For continuity, let $(x_n,z_n)\to(x,z)$ and $\pi_n=\proj_{P_b(x_n)}(z_n)$.
The $\pi_n$ lie in the private box. Let $\bar\pi$ be the limit of a
subsequence. Passing to the limit in $\widehat A_b\pi_n\le\widehat a_b+
\widehat E_bx_n$ shows $\bar\pi\in P_b(x)$. For any $\zeta\in P_b(x)$, the
Hausdorff bound gives $\zeta_n\in P_b(x_n)$ with $\zeta_n\to\zeta$, and
$\norm{\pi_n-z_n}\le\norm{\zeta_n-z_n}$ gives $\norm{\bar\pi-z}\le
\norm{\zeta-z}$ in the limit. So $\bar\pi=\proj_{P_b(x)}(z)$. Since every
subsequence has a further subsequence with this limit, $\pi_n$ converges to
it.
\end{proof}

\begin{lemma}[Integrated displacement]\label{lem:rec-displacement}
Let $m\ge2$ with $v_b(m-1)\le r-1$. For every $u$ with $h_b(u)>0$,
\begin{equation}\label{eq:rec-cs-displacement}
  h_b(u)\,\norm{\bar x_b(u)-u}_2^2\le\sum_{i\in B_b}
  L_b\bigl(K_{m,b}(\cdot,u)\,(x_i-u_i)^2\bigr),
\end{equation}
and, defining the left side to be zero where $h_b(u)=0$,
\[
  \int h_b(u)\norm{\bar x_b(u)-u}_2^2\,d\arcs^{B_b}(u)\le v_bV_m,\qquad
  \int h_b(u)\norm{\bar x_b(u)-u}_2\,d\arcs^{B_b}(u)\le\sqrt{v_bV_m}.
\]
\end{lemma}

\begin{proof}
Write $K=K_{m,b}(\cdot,u)$, $h=h_b(u)$, and $U_i=L_b(Kx_i)$. Then
$L_b(K(x_i-u_i)^2)-h(\bar x_{b,i}-u_i)^2=L_b(Kx_i^2)-U_i^2/h\ge0$, by the
$2\times2$ minor of $\widehat M_b(u)\succeq0$ formed by the entries for $1$ and
$x_i$ (\cref{lem:conditional-matrix}(c)). Summing over $i$ gives
\eqref{eq:rec-cs-displacement}, which also holds where $h=0$. The
right side is a polynomial in $(x,u)$ inside $L_b$. Integrating over $u$
removes the factors $\Kj(x_j,u_j)$, $j\ne i$, and \cref{lem:rec-kernel}(b)
gives
\[
  \int L_b\bigl(K(x_i-u_i)^2\bigr)\,d\arcs^{B_b}(u)
  =V_m+\frac{3-2m}{a_0}L_b(x_i^2)\le V_m,
\]
since $L_b(x_i^2)\ge0$. Sum over $i$. The last bound follows from the
Cauchy--Schwarz inequality and $\int h_b\,d\arcs^{B_b}=1$.
\end{proof}

The constant $V_m$ is the best possible in this lemma: for the evaluation
functional at $x=0$, equality holds in \eqref{eq:rec-cs-displacement} and
in the integral. This sharpness concerns only the displacement estimate.

\begin{theorem}[Affine recourse]\label{thm:affine-recourse}
Consider the recourse model with complete recourse and private convexity.
Let $r\ge\max(w+1,d)$ and $m=\lceil r/w\rceil$. Every feasible family
$(L_b)$ of order $r$ admits a probability law $\nu$ on the feasible set with
\begin{equation}\label{eq:rec-affine-rounding}
  \int f\,d\nu\le\sum_bL_b(f_b)+\frac{3\Abud}{\Dm}+\sqrt{V_m}\,\mathcal R .
\end{equation}
Consequently
\[
  0\le\fstar-\rhorec_r=\fstar-\lambda^{\mathrm{rec}}_r
  \le\frac{3\Abud}{2m^2+1}+\sqrt{V_m}\,\mathcal R
  \le\frac{3w^2\Abud}{2r^2}+\frac{\sqrt3\,w\,\mathcal R}{r}.
\]
\end{theorem}

\begin{proof}
Since $m=\lfloor(r-1)/w\rfloor+1\ge2$, we have $v_b(m-1)\le r-1$, so
\cref{lem:conditional-matrix} applies with $\tau=1$ and
\cref{lem:rec-displacement} applies. Define
$\phi_b(u)=\proj_{P_b(u)}(\bar y_b(u))$ if $h_b(u)>0$ and
$\phi_b(u)=\proj_{P_b(u)}(0)$ otherwise. By \cref{lem:rec-hoffman}, $\phi_b$
is continuous on the open set $\{h_b>0\}$ and on its closed complement, so
it is Borel, and $\phi_b(u)\in P_b(u)$ for every $u$; complete recourse is
used here at every output point.

Fix $u$ with $h_b(u)>0$. By \eqref{eq:rec-source-feasible} and
\eqref{eq:rec-repair}, $\norm{\phi_b(u)-\bar y_b(u)}_2\le\kappa_b\norm{E_b}_2
\norm{\bar x_b(u)-u}_2$. Both points lie in the private box, hence so does
the segment between them, and
\[
  f_b(u,\phi_b(u))\le f_b(u,\bar y_b(u))+\Gamma_b\norm{\bar x_b(u)-u}_2 .
\]
Multiply by $h_b(u)$ and use \eqref{eq:rec-jensen}, whose proof uses only
$M_b(u)\succeq0$ and $Q_b(u)\succeq0$ and not feasibility of $\bar y_b(u)$ at
$u$. Together with the case $h_b(u)=0$, where both sides vanish, this gives
for every $u$
\[
  h_b(u)f_b(u,\phi_b(u))\le\ip{H_b(u)}{M_b(u)}
  +\Gamma_b\,h_b(u)\norm{\bar x_b(u)-u}_2 .
\]
Integrate, and use \cref{lem:rec-matrix-cost,lem:rec-displacement}:
$\int h_bf_b(u,\phi_b(u))\,d\arcs^{B_b}\le L_b(f_b)+3\Abud_b/\Dm
+\Gamma_b\sqrt{v_bV_m}$. Summing over bags and applying \cref{lem:rec-glue}
gives \eqref{eq:rec-affine-rounding}. The infimum over feasible families and
\cref{prop:rec-dual} give the gap bound. Finally $m\ge r/w$, so
$\Dm>2r^2/w^2$ and $\sqrt{V_m}<\sqrt{6/\Dm}<\sqrt3\,w/r$.
\end{proof}

If every $E_b$ is zero, then $\mathcal R=0$ and \cref{thm:fixed-recourse}
already gives the second-order bound without the one-degree reserve. The
objective smoothing still costs only $O(r^{-2})$; the first-order term comes
entirely from moving the private mean from the fiber at $\bar x_b(u)$ to the
fiber at $u$. Repairing a relaxation's private first moments by a
fixed-matrix error bound has precedents in moment relaxations for
distributionally robust optimization~\cite{zhang2025-wasserstein-moment}, and
Hoffman's error bound is classical~\cite{hoffman1952-approximate-solutions},
with computational treatments of its constants~\cite{pena2018-hoffman-constants}.
The contribution here is the combination with exactly consistent kernel
densities in this sparse hierarchy and the resulting explicit rate.

The two structural hypotheses are used essentially. Complete recourse is
needed because the kernel can place mass at every shared point, and a
point with an empty fiber admits no repair. The fixed matrix $A_b$ is
needed twice: to pass from (R2) to source feasibility by taking
conditional means, and for the uniform Hoffman bound. With a
shared-dependent matrix neither step is available. For example
$P(x)=\{y\in[-1,1]:xy=0\}$ has nonempty fibers, but $P(0)=[-1,1]$ and
$P(x)=\{0\}$ for $x\ne0$, so the fibers are not Hausdorff continuous at $0$.
Discontinuous recourse of this type is a known
phenomenon~\cite{zhong2024-two-stage-polynomial}. The constants $\kappa_b$ can
be large for nearly inconsistent constraint systems, and the theorem gives
no method for computing a sharp Hoffman constant.

\subsection{The inverse-order rate is sharp}\label{sec:rec-sharp}

We now show that the order $1/r$ in \cref{thm:affine-recourse} cannot be
improved under its hypotheses. The example has one shared coordinate $x$,
a private $y_1\in[-1,1]$ in the first bag, a private $y_2\in[0,1]$ in the
second bag, and
\begin{equation}\label{eq:rec-sharp-problem}
  \min\ f=-xy_1+y_2\quad\text{subject to}\quad
  x,y_1\in[-1,1],\quad y_2\in[0,1],\quad y_2\ge x,\quad y_2\ge-x .
\end{equation}
For fixed $x$, $\min_{y_1}(-xy_1)=-\abs{x}$ and the smallest feasible $y_2$
is $\abs{x}$, so $\fstar=0$. Both private problems are linear programs.

We consider two different hierarchies for this problem.

\emph{Private-degree-two hierarchy.} Substitute $y_2=(1+\zeta)/2$ with
$\zeta\in[-1,1]$. The bags are $B_1=B_2=\{1\}$ with private vectors $y_1$ and
$\zeta$, the first fiber is $[-1,1]$ (no rows), and the second fiber has
the rows $1-2x+\zeta\ge0$ and $1+2x+\zeta\ge0$, that is, $A_2=(-1,-1)^\top$,
$a_2=(1,1)^\top$, and $E_2=(-2,2)^\top$. The objectives are $f_1=-xy_1$ and
$f_2=(1+\zeta)/2$. This is a recourse model with complete recourse ($\zeta=1$
is always feasible) and private convexity, and \cref{def:rec-hierarchy}
gives values $\rhorec_r$.

\emph{Total-degree hierarchy.} Use the bags $\{y_1,x\}$ and $\{x,y_2\}$ and the
generator lists $G_1=(1-y_1^2,\,1-x^2)$ and
$G_2=(1-x^2,\,y_2,\,1-y_2,\,y_2-x,\,y_2+x)$. Let $T_{b,r}$ be the full
preordering of $G_b$ truncated at total degree $2r$: finite sums of
$\sigma_J\prod_{g\in J}g$ over subsets $J\subseteq G_b$, with $\sigma_J$ a sum
of squares and each summand of total degree at most $2r$. Let
$\rho^{\mathrm{tot}}_r$ be the infimum of $-L_1(xy_1)+L_2(y_2)$ over linear
functionals $L_b$ on polynomials of total degree at most $2r$ in the
variables of bag $b$ with $L_b(1)=1$, $L_b\ge0$ on $T_{b,r}$, and
$L_1(x^j)=L_2(x^j)$ for $0\le j\le2r$. Let $\lambda^{\mathrm{tot}}_r=\sup
\{\lambda:f-\lambda\in T_{1,r}+T_{2,r}\}$. As in \cref{lem:weak-duality},
\cref{lem:zero-decomposition} for the junction tree with these two bags
gives $\lambda^{\mathrm{tot}}_r\le\rho^{\mathrm{tot}}_r\le\fstar=0$.

The two hierarchies are not compared by inclusion: the second has
products of private constraints and higher private degrees, while the
first has mixed moments of joint degree up to $2r+2$.

\begin{lemma}[A moment-matching witness for $\abs{x}$]\label{lem:rec-abs-witness}
For every integer $n\ge0$ there are probability measures $\alpha,\beta$ on
$[-1,1]$ with $\int x^j\,d\alpha=\int x^j\,d\beta$ for $0\le j\le n$ and
\[
  \int\abs{x}\,d\alpha-\int\abs{x}\,d\beta\ge\frac2{9\pi(n+2)} .
\]
\end{lemma}

\begin{proof}
Let $N$ be the even integer with $n+1\le N\le n+2$ and let
$F_N(s)=\sum_{\abs{j}<N}(1-\abs{j}/N)e^{ijs}$ be the Fejér kernel, which is
nonnegative with $\int F_N\,ds/(2\pi)=1$. Put
\[
  Q(\theta)=-\cos\bigl(2N(\theta-\tfrac\pi2)\bigr)F_N\bigl(\theta-\tfrac\pi2\bigr)
  =-\sum_{\abs{j}<N}\Bigl(1-\frac{\abs{j}}N\Bigr)\cos\bigl((2N+j)(\theta-\tfrac\pi2)\bigr),
\]
where the second form follows from $\cos(2Ns)e^{ijs}+\cos(2Ns)e^{-ijs}$
pairing. Then $\int\abs{Q}\,d\theta/(2\pi)\le1$, and all frequencies $2N+j$
lie in $[N+1,3N-1]$. Since $p(\cos\theta)$ is a trigonometric polynomial of
frequencies at most $n<N+1$ for $\deg p\le n$, $\int p(\cos\theta)Q(\theta)\,
d\theta/(2\pi)=0$.

The function $\abs{\cos\theta}$ is even with period $\pi$; its cosine
coefficients $c_k$ vanish for odd $k$, and for even $k=2l\ge2$ a direct
integration gives $c_k=-4(-1)^l/(\pi(k^2-1))$. Since
$\cos(k(\theta-\tfrac\pi2))=\cos(k\theta)\cos(k\pi/2)+\sin(k\theta)\sin(k\pi/2)$
and $\abs{\cos\theta}$ is even,
\[
  \int\abs{\cos\theta}\cos\bigl(k(\theta-\tfrac\pi2)\bigr)\frac{d\theta}{2\pi}
  =\frac{c_k\cos(k\pi/2)}2=-\frac2{\pi(k^2-1)}\quad(k\ge2\text{ even}),
\]
and it is zero for odd $k$. The frequency $2N+j$ is even exactly when $j$
is even, and then $(2N+j)^2-1<9N^2$. Hence
\[
  \int\abs{\cos\theta}\,Q(\theta)\frac{d\theta}{2\pi}
  =\sum_{\substack{\abs{j}<N\\ j\text{ even}}}\Bigl(1-\frac{\abs{j}}N\Bigr)
  \frac2{\pi((2N+j)^2-1)}
  \ge\frac2{9\pi N^2}\cdot\frac N2=\frac1{9\pi N},
\]
because, with $N=2M$, the even-index weights sum to
$1+2\sum_{i=1}^{M-1}(1-i/M)=M$.

Let $\varsigma$ be the image of the signed measure $Q(\theta)\,d\theta/(2\pi)$
under $\theta\mapsto\cos\theta$. It has total variation at most one, total
mass zero (frequency zero is absent), annihilates all polynomials of degree
at most $n$, and satisfies $\int\abs{x}\,d\varsigma\ge1/(9\pi N)>0$. Write its
Jordan decomposition $\varsigma=\varsigma_+-\varsigma_-$. Both parts have the
same mass $t$, with $0<t\le\frac12$. The probability measures
$\alpha=\varsigma_+/t$ and $\beta=\varsigma_-/t$ have equal moments through
degree $n$, and $\int\abs{x}\,d(\alpha-\beta)=t^{-1}\int\abs{x}\,d\varsigma\ge
2/(9\pi N)\ge2/(9\pi(n+2))$.
\end{proof}

The lemma is a quantitative form of the classical fact that $\abs{x}$ is
approximated by polynomials of degree $n$ only to accuracy of order $1/n$;
sharper constants are classical. The identity relating the best moment
matching difference to best uniform approximation is recorded in
\cref{lem:moment-matching-general}; we need only the explicit lower bound.

\begin{theorem}[Sharpness of the inverse-order rate]\label{thm:affine-sharp}
For problem \eqref{eq:rec-sharp-problem} and every integer $r\ge2$:
\begin{enumerate}[label=(\alph*)]
\item \emph{(private-degree-two hierarchy)}
  \[
    \frac1{9\pi(r+1)}\le-\rhorec_r\le\frac3{2r^2+1}+\sqrt{2V_r}
    <\frac3{2r^2+1}+\frac{2\sqrt3}{\sqrt{2r^2+1}} ;
  \]
\item \emph{(total-degree hierarchy)}
  \[
    \frac1{9\pi(r+1)}\le-\rho^{\mathrm{tot}}_r\le-\lambda^{\mathrm{tot}}_r
    \le2\sqrt{V_r}<\frac{2\sqrt6}{\sqrt{2r^2+1}}\le\frac{2\sqrt3}r .
  \]
\end{enumerate}
In both hierarchies the gap has exact order $1/r$, and neither is exact at
any finite order.
\end{theorem}

\begin{proof}
\emph{Lower bounds.} Apply \cref{lem:rec-abs-witness} with $n=2r$. Give the
first bag the law of $(x,\operatorname{sign}(x))$ with $x\sim\alpha$ and
$\operatorname{sign}(0)=0$, and the second bag the law of $(x,\abs{x})$ with
$x\sim\beta$ (in the private-degree-two encoding, $\zeta=2\abs{x}-1$). These
laws are supported on the local feasible sets. Integration against them
defines functionals that are nonnegative on every polynomial tested in
either hierarchy, since each such polynomial is nonnegative on the local
feasible set, and the separator moments agree through degree $2r$. Their
objective is $-\int\abs{x}\,d\alpha+\int\abs{x}\,d\beta\le-2/(9\pi(2r+2))$.
Hence both moment values are at most $-1/(9\pi(r+1))$.

\emph{Upper bound in (a).} We apply \cref{thm:affine-recourse} with $w=1$,
$d=1$, and $m=r$. Here $H_1(x)=xH_{1,(1)}$ with
$H_{1,(1)}=\bigl(\begin{smallmatrix}0&-1/2\\-1/2&0\end{smallmatrix}\bigr)$,
so $\Abud_1=1$, and $H_2$ is constant, so $\Abud=1$. The first bag has no
rows, so $\Gamma_1=0$. For the second, $\Lambda_2=\tfrac12$ and
$\norm{E_2}_2=2\sqrt2$. All rows of $\widehat A_2=(-1,-1,1,-1)^\top$ have
absolute value one, so a nonempty solution set is an interval cut out by
half-lines, and the distance from a point to that interval is the largest
distance to one of the half-lines, which equals the largest positive
residual. Hence $\kappa_2=1$ is admissible, $\Gamma_2=\sqrt2$, and
$\mathcal R=\sqrt2$. The bound follows from $V_r<6/(2r^2+1)$.

\emph{Upper bound in (b).} Let $m=r$, let $s_m(x)=\int\Kj(x,t)\operatorname{sign}(t)\,
d\arcs(t)$ and $\delta_m=2\sqrt{V_m}$. Since $\Kj\ge0$ integrates to one,
$\abs{s_m}\le1$. By \eqref{eq:rec-circle-kernel},
$s_m=\sum_{k\ge1}2\kmult_kT_k\int T_k\operatorname{sign}\,d\arcs$, and the integrals
vanish for even $k$; so $s_m$ is odd of degree at most $2m-3$. For all
$x,t\in[-1,1]$, $0\le\abs{x}-x\operatorname{sign}(t)\le2\abs{t-x}$: the middle
expression is $0$ if the signs agree, $\abs{x}$ if $t=0$, and $2\abs{x}\le
2\abs{t-x}$ if the signs are opposite. Integrating against
$\Kj(x,t)\,d\arcs(t)$ and using the Cauchy--Schwarz inequality and
\eqref{eq:rec-displacement-identity},
\[
  0\le\abs{x}-xs_m(x)\le2\Bigl(\int\Kj(x,t)(x-t)^2\,d\arcs(t)\Bigr)^{1/2}
  \le\delta_m .
\]
Put $p_m=xs_m+\delta_m$, so $p_m\ge\abs{x}$ on $[-1,1]$ and
$\deg p_m\le2m-2$. Then
\[
  f+\delta_m=\underbrace{\Bigl(\frac{1+y_1}2(p_m-x)+\frac{1-y_1}2(p_m+x)\Bigr)}_{=\,p_m-xy_1}
  +\underbrace{\Bigl(\frac{1+s_m}2(y_2-x)+\frac{1-s_m}2(y_2+x)\Bigr)}_{=\,y_2-xs_m}.
\]
The polynomials $p_m\pm x$ and $(1\pm s_m)/2$ are nonnegative on $[-1,1]$
and have degree at most $2m-2$. By \cref{lem:interval-sos} each equals
$\sigma_0+(1-x^2)\sigma_1$ with $\sigma_0\in\sos_{2m-2}$ and
$\sigma_1\in\sos_{2m-4}$. Also $(1\pm y_1)/2=(1\pm y_1)^2/4+(1-y_1^2)/4$.
Expanding, the first bracket is a sum of terms $\sigma$, $(1-x^2)\sigma$,
$(1-y_1^2)\sigma$, and $(1-y_1^2)(1-x^2)\sigma$ with sums of squares $\sigma$,
each of total degree at most $2m$; it lies in $T_{1,m}$. The second bracket
is a sum of terms $\sigma(y_2\mp x)$ and $(1-x^2)\sigma(y_2\mp x)$ of total
degree at most $2m-1$; it lies in $T_{2,m}$. Thus
$\lambda^{\mathrm{tot}}_r\ge-\delta_r$. The final inequalities use
$V_r<6/(2r^2+1)$ and $6/(2r^2+1)\le3/r^2$.
\end{proof}

The lower bound is a statement about actual local measures: it holds for
every relaxation whose local conditions are valid for measures on the local
feasible sets and whose separator information is polynomial of degree at
most $2r$, however strong its local positivity conditions are. The upper
bounds are proved separately for each hierarchy; the certificate in (b)
uses the products $(1-x^2)(y_2\mp x)$ of the full preordering. Qualitative
failure of finite sparse exactness with a dense certificate is known from a
two-bag example whose separator function is
rational~\cite{nie2026-sparse-tightness}; here the separator function is
$\abs{x}$, and the rate is identified. Bernstein's theorem on best uniform
approximation of $\abs{x}$ gives the exact asymptotics of the local-measure
gap; the leading constants of the finite hierarchy gaps are not determined
here.

\begin{remark}[Merging or splitting removes the gap for this instance]
\label{rem:merge-split}
The obstruction concerns the fixed two-bag decomposition, not the
difficulty of \eqref{eq:rec-sharp-problem}. The identity
\[
  f=\frac{(1+y_1)^2}4(y_2-x)+\frac{(1-y_1)^2}4(y_2+x)+\frac{1-y_1^2}2\,y_2
\]
is a degree-three certificate in the dense preordering of the merged bag
$\{y_1,x,y_2\}$; its last term multiplies generators of different bags.
Alternatively, splitting the problem at $x=0$ and adding the generator $x$
(respectively $-x$) to the first bag gives the exact certificates
$f=(y_2-x)+x(1-y_1)$ and $f=(y_2+x)+(-x)(1+y_1)$, where
$1\mp y_1=\tfrac12(1\mp y_1)^2+\tfrac12(1-y_1^2)$; each has degree three and
keeps the two bags. These are properties of this instance only.
\end{remark}

\subsection{Comparison with discretization}\label{sec:rec-grid}

\begin{corollary}[Comparison of a fixed grid with every moment family]
\label{cor:rec-grid}
Let $d_\infty$ be the largest individual shared degree in any entry of
$H_b$, with zero for constant entries. Under the hypotheses and parameter
choice of \cref{thm:fixed-recourse}, put
$N=m+\lfloor d_\infty/2\rfloor$. Under those of
\cref{thm:affine-recourse}, instead put
$N=m+\lfloor\max(d_\infty,2)/2\rfloor$. In both cases let
$\Xi_N=\{\cos((2j-1)\pi/(2N)):1\le j\le N\}$ and
$F_b(u)=\min_{y\in P_b(u)}f_b(u,y)$. Every feasible family satisfies
\[
 \min_{x\in\Xi_N^n}\sum_bF_b(x_{B_b})
 \le\sum_bL_b(f_b)+\frac{3\Abud}{\Dm}
 +\begin{cases}0,&\text{fixed private domain},\\
               \sqrt{V_m}\,\mathcal R,&\text{affine recourse}.
   \end{cases}
\]
After the $\sum_bN^{v_b}$ local QP values have been computed, junction-tree
dynamic programming finds the grid minimum in $O(tN^w)$
arithmetic and comparison operations. This count excludes construction
and certification of the QP tables.
\end{corollary}

\begin{proof}
By \cref{lem:ker-quadrature}, equal-weight quadrature on $\Xi_N$ integrates
univariate polynomials through degree $2N-1$ exactly. The masses
$N^{-v_b}h_b(u)$ at tensor grid nodes are nonnegative, sum to one, and have
exactly matching separator marginals: quadrature deletes each
nonseparator kernel factor, and the remaining separator density is the
one in \cref{lem:rec-glue}. Finite-state gluing therefore gives a global
shared grid law with these bag marginals.

In the fixed case, the pointwise Jensen inequality gives
$h_b(u)F_b(u)\le\ip{H_b(u)}{M_b(u)}$, including zero-density nodes. The
matrix expression is polynomial of individual degree at most
$2m-2+d_\infty\le2N-1$. Exact quadrature and
\cref{lem:rec-matrix-cost} bound its grid expectation by
$L_b(f_b)+3\Abud_b/\Dm$.

In the affine case, use the repaired decision from
\cref{thm:affine-recourse} at each grid node. The pointwise matrix and
repair bounds give
$h_bF_b\le\ip{H_b}{M_b}+\Gamma_bh_b\norm{\bar x_b-u}_2$.
The displacement polynomials
$L_b(K_{m,b}(\cdot,u)(x_i-u_i)^2)$ have individual output degree at most
$2m$, and the chosen $N$ satisfies $2N-1\ge2m$. Hence quadrature
reproduces the displacement integral in \cref{lem:rec-displacement}.
Finite-sum Cauchy--Schwarz bounds the grid expectation of
$h_b\norm{\bar x_b-u}_2$ by $\sqrt{v_bV_m}$. Summing the bag bounds
gives the stated cost expectation, and some global grid assignment has
no larger cost. The dynamic program visits at most $N^{v_b}$ entries per
bag and per parent--child update; the tree has $t-1$ edges, giving
$O(tN^w)$ once the local tables are supplied. Grid nodes are algebraic
and coefficients may be real; the count is an exact-arithmetic or oracle
statement, not a bit-complexity bound.
\end{proof}

\begin{remark}[Grid baseline]\label{rem:grid-baseline}
The exponents of \cref{thm:fixed-recourse,thm:affine-recourse} are also
reached without any semidefinite program, by discretizing the shared
coordinates and solving convex quadratic programs. Let $N\ge1$,
$\xi_j=\cos((2j-1)\pi/(2N))$ for $1\le j\le N$, $\Xi=\{\xi_1,\dots,\xi_N\}$,
and $F_b(x)=\min_{y\in P_b(x)}f_b(x,y)$. Then
\begin{equation}\label{eq:rec-grid}
  \min_{x\in\Xi^n}\sum_bF_b(x_{B_b})\le\fstar+\frac{\pi^2\Abud}{8N^2}
  +\begin{cases}
      0,&\text{fixed private domain},\\
      \pi\mathcal R/N,&\text{complete affine recourse}.
    \end{cases}
\end{equation}
Once the $\sum_bN^{v_b}$ values $F_b$ on the grid are known, ordinary
junction-tree dynamic programming minimizes the left side.

To prove \eqref{eq:rec-grid}, let $(x^*,y^*)$ be optimal and write
$x^*_i=\cos\theta^*_i$ with $\theta^*_i\in[0,\pi]$. The angles
$(2j-1)\pi/(2N)$, $j\in\Z$, have spacing $\Delta=\pi/N$, and all their
cosines lie in $\Xi$. Round each shared coordinate independently: pick the
two neighboring angles of $\theta^*_i$ with the probabilities that make the
mean angle equal to $\theta^*_i$, and let $X_i$ be the cosine of the chosen
angle. Linear interpolation of $\theta\mapsto\cos(k\theta)$ has error at most
$k^2\Delta^2/8$, so $\abs{\mathbb E\,T_k(X_i)-T_k(x^*_i)}\le k^2\Delta^2/8$.
By independence and since all factors have absolute value at most one,
$\abs{\mathbb E\,T_\alpha(X)-T_\alpha(x^*)}\le\Delta^2\sum_i\alpha_i^2/8$.
Since the entries of $\zy^*_b$ have absolute value at most one,
$\mathbb E\sum_bf_b(X_{B_b},y^*_b)\le\fstar+\pi^2\Abud/(8N^2)$. With a fixed
private domain $y^*_b\in P_b$, and $F_b(X_{B_b})\le f_b(X_{B_b},y^*_b)$. Under
complete recourse, replace $y^*_b$ by its projection onto $P_b(X_{B_b})$.
Since $\abs{X_i-x^*_i}\le\Delta$, \eqref{eq:rec-repair} and the Lipschitz
bound $\Lambda_b$ increase the cost of bag $b$ by at most
$\Gamma_b\sqrt{v_b}\,\pi/N$. Some realization of $X$ does no worse than
the expectation.

With $N$ proportional to $r$, \eqref{eq:rec-grid} has the same exponents as
the two theorems; the argument does not even use private convexity, which
only makes the local problems convex. Chebyshev-grid bounds of this type
are established~\cite{piazzon2018-chebyshev-grids}. The theorems of this
section are of a different kind: they bound every feasible point of the
private-degree-two hierarchy, and hence its moment value and certificate
value, rather than the value of a particular discrete search. We make no
comparison of the computational effort of the two approaches.
\end{remark}
